% Part: incompleteness
% Chapter: representability-in-q
% Section: representing-relations

\documentclass[../../../include/open-logic-section]{subfiles}

\begin{document}

\olfileid[fr]{inc}{req}{rel}

\olsection{Représentation des relations}

Précisons ce que signifie représenter une \emph{relation}.

\begin{defn}
\ollabel{defn:representing-relations}
Une relation $R(x_0,\dots,x_k)$ sur les entiers naturels est
{\em représentable dans $\Th{Q}$} s'il existe une
!!{formula}~$!A_R(x_0,\dots,x_k)$ telle que, pour tous
$n_0,\dots,n_k$, $\Th{Q}$ démontre
$!A_R(\num{n_0},\dots,\num{n_k})$ si
$R(n_0,\dots,n_k)$ est vraie, et démontre
$\lnot !A_R(\num{n_0},\dots,\num{n_k})$ si elle est fausse.
\end{defn}

\begin{thm}
\ollabel{thm:representing-rels}
Une relation est représentable dans~$\Th{Q}$ si et seulement
si elle est calculable.
\end{thm}

\begin{proof}
Pour le sens direct, supposons que la
!!{formula}~$!A_R(x_0,\dots,x_k)$ représente
$R(x_0,\dots,x_k)$. Voici un algorithme qui décide~$R$ :
sur l'entrée $n_0,\dots,n_k$, recherchons en parallèle une
preuve de $!A_R(\num{n_0},\dots,\num{n_k})$ et une preuve
de $\lnot !A_R(\num{n_0},\dots,\num{n_k})$.
L'hypothèse de représentabilité garantit que l'une des deux
recherches aboutit. Si la première aboutit, répondons « oui »;
si c'est la seconde, répondons « non ». La cohérence de
$\Th{Q}$ exclut que les deux réponses soient démontrables.

Réciproquement, supposons $R(x_0,\dots,x_k)$ calculable.
Par définition, sa fonction caractéristique
$\Char{R}(x_0,\dots,x_k)$ est calculable. D'après le
\olref[int]{thm:representable-iff-comp}, elle est
représentée par une !!{formula}, notée
$!A_{\Char{R}}(x_0,\dots,x_k,y)$. Prenons pour
$!A_R(x_0,\dots,x_k)$ la !!{formula}
$!A_{\Char{R}}(x_0,\dots,x_k,\num{1})$.

Si $R(n_0,\dots,n_k)$ est vraie, alors
$\Char{R}(n_0,\dots,n_k)=1$. La théorie $\Th{Q}$ démontre
donc $!A_{\Char{R}}(\num{n_0},\dots,\num{n_k},\num{1})$,
c'est-à-dire $!A_R(\num{n_0},\dots,\num{n_k})$.
Si, au contraire, $R(n_0,\dots,n_k)$ est fausse, alors
$\Char{R}(n_0,\dots,n_k)=0$. La seconde condition de
représentabilité de cette fonction donne
\[
\lforall[y][(!A_{\Char{R}}(\num{n_0}, \dots, \num{n_k}, y) \lif y =
  \num{0})].
\]
Puisque $\Th{Q}$ démontre aussi
$\eq/[\num{0}][\num{1}]$, elle démontre
$\lnot !A_{\Char{R}}(\num{n_0},\dots,\num{n_k},\num{1})$,
donc $\lnot !A_R(\num{n_0},\dots,\num{n_k})$.
\end{proof}

\begin{prob}
Montrez que, si $R$ est représentable dans~$\Th{Q}$,
sa fonction caractéristique $\Char{R}$ l'est aussi.
\end{prob}

\end{document}
